\subsection{Feature Ablation Analysis}
\label{sec:feature_ablation}

The quantification of classification accuracy changes from single-domain schemes toward dual- and tri-domain configurations is summarized in Table~\ref{tab:VII}, Table~\ref{tab:VIII}, Table~\ref{tab:IX}, and Table~\ref{tab:X}. In the SVM model, the representation transition from Face (90.83\%) to Emotion $\oplus$ Face (93.29\%) increases accuracy by 2.46\% (0.0246). The addition of a third domain in the tri-domain Face $\oplus$ Emotion $\oplus$ Age configuration yields a further performance improvement to 93.70\%, which reflects a cumulative increase of 2.87\% (0.0287) over the single Face feature. This pattern of progressive improvement indicates the presence of additional discriminative information from each combined visual representation domain, although the interaction mechanisms among latent feature dimensions cannot be definitively confirmed solely from empirical test results.

An analysis of the contribution of the age domain (Age) representation dissects the role of latent aging information within the demographic classification system. Across all evaluated classifiers, Age constitutes the single-domain configuration with the lowest accuracy attainment, as recorded at 87.64\% for SVM and 69.63\% for GNB. Nevertheless, combining the Age feature with facial biometric features in the dual-domain Face $\oplus$ Age scheme yields an accuracy of 92.55\%, providing an improvement of 1.72\% (0.0172) over the pure Face feature performance (90.83\%) on the SVM model. These empirical results indicate the likely presence of additional discriminative information from age-associated representations that complements the facial morphological profile, without claiming that such informational contribution is orthogonal or strictly complementary theoretically.
